%%%%%%%%%%%%%%%%%%%%%%%%%%%%%%%%%%%%%%%%%%%%%%%%%%%%%%%%%%%%%%
%%%%%%%%%%%% CA-0721 Probabilidad %%%%%%%%%%%%%%%%%%%%%%%%%%%%%
%%%%%%%%%%%%%%%%%%%%%%%%%%%%%%%%%%%%%%%%%%%%%%%%%%%%%%%%%%%%%%
\newpage
\subsection*{CA-0721 Probabilidad}
\addcontentsline{toc}{subsection}{CA-0721 Probabilidad}

\hypertarget{CA0721}{}
\begin{refsection}
\begin{table}[H]
\centering{
\begin{tabular}{|ll|}
\hline
\textbf{Año:} II  & \textbf{Requisitos:} CA-0351, CA-0352 \\ 
\textbf{Ciclo:} II  & \textbf{Correquisitos:} Ninguno \\ 
\textbf{Tipo de curso:} Teórico & \textbf{Horas semanales:} 5 \\ 
\textbf{Créditos:} 4 & \textbf{Virtualidad:} PP, P \\ \hline
\end{tabular}
}
\end{table}

\subsubsection*{Descripción del curso}

Este curso está dirigido a estudiantes que completaron los cursos de cálculo en múltiples variables y análisis matemático por lo que se ubica en el segundo año y segundo ciclo. Este curso conecta el área de matemáticas de la carrera con las demás áreas, en particular de estadística y de Ciencias Actuariales, por lo que es medular y requisito para los siguientes cursos. El propósito del curso es que el/la estudiante conozca sobre los principales resultados de la teoría de probabilidades y cómo se aplican a problemas de cálculo de probabilidades.

\subsubsection*{Objetivo general}

Introducir los resultados teóricos y propiedades principales de la teoría de probabilidades y su aplicación en la solución de problemas relacionados.

\subsubsection*{Objetivos específicos}

Al finalizar el curso el/la estudiante estará en la capacidad de:

\begin{enumerate}
\item Calcular probabilidades que se derivan de fenómenos aleatorios simples. (SC02.)
\item Aplicar propiedades básicas de las variables y vectores aleatorios en la solución de distintos problemas. (SC02.)
\item Utilizar propiedades de transformadas para el estudio de variables y vectores aleatorios. (SC02.)
\item Interpretar la noción de independencia y esperanza condicional para variables aleatorias. (SC02.)
\item Identificar diferentes modos de convergencia de las variables aleatorias, tanto en aplicaciones sencillas como en leyes de grandes números y teoremas de límite central. (SC02.)
\end{enumerate}

\subsubsection*{Contenidos}

\begin{enumerate}
\item \textbf{Espacios de probabilidad}
\begin{enumerate}
\item Experimentos aleatorios.
\item Espacio muestral.
\item $\sigma$-álgebra de eventos.
\item Medidas de probabilidad.
\item Cálculo de probabilidades.
\item Probabilidad geométrica.
\item Modelo de probabilidad de Laplace.
\item Probabilidad condicional.
\item Independencia de eventos.
\end{enumerate}
\item \textbf{Variables aleatorias}
\begin{enumerate}
\item Variables aleatorias y operaciones sobre ellas.
\item Función de distribución.
\item Variables aleatorias discretas y absolutamente continuas.
\item Función de masa y densidad.
\item Esperanza, varianza y momentos.
\item Fórmula de cambio de variable.
\item Función generadora de momentos.
\item Función característica.
\end{enumerate}
\item \textbf{Vectores aleatorios}
\begin{enumerate}
\item Función de distribución conjunta.
\item Función de masa y de densidad conjunta.
\item Independencia.
\item Matriz de covarianza.
\item Fórmula del cambio de variable.
\item Función generadora de momentos.
\item Función característica.
\item Distribución normal multivariada.
\end{enumerate}
\item \textbf{Esperanza condicional}
\begin{enumerate}
\item Distribución condicional.
\item Función de masa y densidad condicional.
\item Esperanza condicional elemental.
\item Esperanza condicional a una $\sigma$-álgebra.
\end{enumerate}
\item \textbf{Convergencia de variables aleatorias}
\begin{enumerate}
\item Convergencia casi segura.
\item Convergencia en probabilidad.
\item Convergencia en $L^p$.
\item Convergencia en distribución.
\end{enumerate}
\item \textbf{Leyes de grandes números}
\begin{enumerate}
\item Ley débil de grandes números.
\item Ley fuerte de grandes números.
\end{enumerate}
\item \textbf{Teoremas de límite central}
\begin{enumerate}
\item Teorema de límite central De Moivre--Laplace.
\item Teorema de límite central.
\end{enumerate}
\end{enumerate}

\subsubsection*{Metodología}

El nivel de virtualidad que se empleará en el curso será medio-virtual. Específicamente, ciertos contenidos pueden integrarse dentro de la dinámica virtual, asincrónica y a partir de estrategias de aprendizaje tales como guías de trabajo que buscará el desarrollo de habilidades específicas. También se apoyará en clases magistrales, en donde se dará un énfasis de índole operacional, sin dejar de lado totalmente la rigurosidad matemática. Además se contemplará aplicaciones en Ciencias Actuariales. Finalmente se fomentará el trabajo colaborativo a partir de pequeños proyectos en la línea de una investigación básica e inicial.

\subsubsection*{Evaluación}

\begin{enumerate}
\item Guías de trabajo (pequeños proyectos de investigación) (carácter formativo y sumativo).
\item Exámenes parciales (presenciales de carácter sumativo).
\end{enumerate}

\subsubsection*{Bibliografía recomendada}

\begin{enumerate}
\item L. Blanco-Castañeda, V. Arunachalam, and S. Dharmaraja. \emph{Introduction to Probability and Stochastic Processes with Applications.} John Wiley \& Sons Inc. (2012).
\item G. Grimmett and D. Welsh. \emph{Probability. An introduction.} Oxford University Press, second edition (2014).
\item P. Hoel, S. Port and C. Stone. \emph{Introduction to probability theory.} Cengage Learning, first edition (1972).
\end{enumerate}

\subsubsection*{Bibliografía complementaria}

\begin{enumerate}
\item R. Durrett. \emph{Probability. Theory and Examples.} Cambridge Series in Statistical and Probabilistic Mathematics, Cambridge University Press, fourth edition (2010).
\item G. Grimmett and D. Stirzaker. \emph{Probability and random processes.} Oxford University Press, third edition (2001).
\item A. Shiryaev. \emph{Probability-1.} Springer, third edition (2016).
\end{enumerate}

\end{refsection}
